% !TEX root = ../main.tex
\chapter{Newtonian kinematics (Newtonsche Kinematik)}\label{ch:c08}
\begin{paracol}{2}

\begin{sources}
\begin{tabularx}{\linewidth}{@{}l l L@{}}
\toprule
\textbf{Video} & \textbf{Length} & \textbf{Content}\\
\midrule
\seg{8}{1}{Lecture 8: Newtonian kinematics} & 1:46:46 & Part B of the course (Newtonian spacetime and the precise
formulation of the three laws): the definition of a Newtonian spacetime (smooth 4-manifold, torsion-free covariant derivative,
absolute time with the two properties Z1 and Z2) with five remarks (outlook on flatness and gravitation, vectors as spatial,
future- or past-pointing, events, the foliation into spaces, stratified atlases); particles as world lines with mass and
charges; the first law in precise form; the affine parametrisation of world lines (with proof); global inertial systems and their
evaluation in components: the straight-line motion in space.\\
\bottomrule
\end{tabularx}

\smallskip
\textbf{Transcript} \file{materials/transcripts/C08-V01.txt} (German YouTube auto-captions).
\end{sources}

\begin{reading}[Where this lecture sits in the other notes]
\emph{Gravity and Light}: lecture 9, ``Newtonian spacetime is curved'' (the same Newtonian spacetime
$(M,\mathcal O,\mathcal A,\nabla,t)$, there with the curvature of $\nabla$ and gravity as spacetime curvature, which
here is only announced). The continuation of this lecture, forces and the second law, is \cref{ch:c10}.
\end{reading}

% ==================================================================
\section{Plan of part B: kinematics}\label{sec:nk-plan}

Today we start part B of the course, \emph{Newtonian spacetime and the precise formulation of the three laws}. Everybody surely remembers
the first lecture, in which the three Newtonian laws were written down in one of the usual formulations and our main insight was that
almost all the load-bearing notions had no precise definition \ts{8}{1}{0:08}. Today this is repaired: all of them receive a completely
precise definition within our mathematical concept formation, and no philosophical interpretation is needed any more of what the
absolute time or an inertial system ``really'' is \ts{8}{1}{0:28}.

We begin with the so-called \textbf{kinematics} (Kinematik) of a theory, here of mechanics: the whole \emph{setup} of the notions
of the theory (what is the absolute time, what do we mean by ``uniform'', and so on) \emph{plus} the first Newtonian axiom (``law'' and ``axiom'' are used as synonyms)
\ts{8}{1}{0:49}. The next lecture treats the \textbf{dynamics}: the second law, forces, the fields that cause them \ts{8}{1}{1:32}.
Today: kinematics and a complete understanding of the first law \ts{8}{1}{1:52}.

% ==================================================================
\section{Newtonian spacetime}\label{sec:nk-spacetime}

We begin with the notion of a \emph{Newtonian spacetime} (Newtonsche Raumzeit). At first sight it may be surprising that one does not first
define space, then time, and only then, as a synthesis of the two, a concept of spacetime. But this is right: one must first define a
spacetime, and only \emph{in} a spacetime can one speak of an absolute time, and only then of spaces. ``The shoe is the other way round'': spacetime
is the fundamental mathematical object, and the word is \emph{not} to be read, mysteriously, as ``space plus time''; it is a new notion, defined now
\ts{8}{1}{2:19}. The adjective ``Newtonian'' distinguishes it from an Einsteinian spacetime, to which we shall come \ts{8}{1}{3:01}.

\begin{definition}[New concepts: Newtonian spacetime (Newtonsche Raumzeit)]\label{def:nk-spacetime}
A \textbf{Newtonian spacetime} is a quintuple $(M,\mathcal O,\mathcal A,\nabla,t)$ consisting of the following three ingredients
\ts{8}{1}{3:25}.
\begin{enumerate}
  \item A \textbf{smooth manifold} $(M,\mathcal O,\mathcal A)$ \textbf{of dimension 4} (``of course I forgot to say four-dimensional'', added
        a little later in the lecture) \ts{8}{1}{3:46}.
  \item A \textbf{covariant derivative} $\nabla$ (equivalently a linear connection $\Gamma$) which is \textbf{torsion free}
        (torsionsfrei) \ts{8}{1}{4:28}.
  \item An \textbf{absolute time} (absolute Zeit): a smooth function $t\colon M\to\R$ with the two properties
        \begin{itemize}
          \item[\textbf{Z1}] (\emph{foliation} of spacetime into three-dimensional spaces, Bl\"atterung): $\mathrm dt\ne0$ at every point; that is, the
                linear map $\mathrm d_pt\colon T_pM\to\R$ is nowhere the zero map;
          \item[\textbf{Z2}] (\emph{the absolute time flows uniformly}): $\nabla_X\,\mathrm dt=0$ for all vector fields $X$.
        \end{itemize}
\end{enumerate}
\end{definition}
\ts{8}{1}{8:00}.
Let us unpack the pieces. A covariant derivative is equivalent to a linear connection (the connection defines the covariant derivative; the converse,
extracting a connection from a covariant derivative, we have not done, but it is possible) \ts{8}{1}{4:48}.

\paragraph{Torsion freeness.} On an earlier problem sheet one had to show that the combination
\begin{equation}\label{eq:nk-torsion}
  T^a{}_{bc}\;:=\;\Gamma^a{}_{bc}-\Gamma^a{}_{cb}
\end{equation}
(components in a chart $(U,x)$ of a linear connection $\Gamma$, with the two lower indices swapped) are the component functions of a \emph{tensor}:
although the $\Gamma$'s do not transform like a $(1,2)$ tensor, this combination does \ts{8}{1}{5:10}. It is the \textbf{torsion tensor}. Of a tensor one
may demand that it \emph{vanish} in all charts: if the components of a tensor vanish in all charts of an atlas they vanish also in every other compatible
chart, which follows directly from the tensor transformation law (and was an exercise). One could \emph{not} demand that the $\Gamma$'s
themselves be zero, which could be true in one chart and false in another. So we demand that the torsion vanish in every chart $(U,x)\in\mathcal A$; the evident consequence is that
\begin{equation}\label{eq:nk-sym}
  \Gamma^a{}_{bc}=\Gamma^a{}_{cb}\qquad\text{in every chart,}
\end{equation}
i.e.\ the $\Gamma$'s are symmetric in the two lower indices \ts{8}{1}{6:34}. As argued, this is a well-defined requirement if made in every chart
\ts{8}{1}{7:39}. So far we have a smooth manifold with a covariant derivative (hence parallel transport, autoparallels and so on), with the restriction of being a special, torsion-free linear connection
\ts{8}{1}{7:18}.
\begin{coursenote}[Why the combination (\ref{eq:nk-torsion}) is a tensor]
Not worked out in the lecture (it was a problem-sheet exercise): the inhomogeneous term of the transformation law of $\Gamma$ in \cref{ch:c06},
$\partial y^a/\partial x^k\,\partial^2x^k/\partial y^b\partial y^c$, is symmetric in $b,c$ (Schwarz's theorem inside one chart),
so it cancels in the difference \eqref{eq:nk-torsion}, which is left with the pure $(1,2)$ tensor law. (Added; not said in class.)
\end{coursenote}

\paragraph{The absolute time.} What is an absolute time mathematically? Simply a smooth function $t\colon M\to\R$ on the manifold with two properties
\ts{8}{1}{8:26}.

\emph{Z1: foliation of spacetime into three-dimensional spaces.} For a smooth function one can form at every point $p\in M$ the object $\mathrm d_pt$,
a linear map $T_pM\to\R$: as before, it eats a vector $X\in T_pM$ and gives the number $X(t)$, the action of the vector on the function
\ts{8}{1}{9:11}. The condition Z1 is that $\mathrm d_pt$ is \emph{not the zero map} for any $p\in M$. It could happen that a function sends every vector at some point to
zero; we want to exclude this \ts{8}{1}{10:40}.

\emph{Z2: the absolute time flows uniformly.} In the first lecture Schuller made fun of the phrase ``time flows uniformly'' (one second per second?).
It now gets a meaning \ts{8}{1}{11:26}. The object $\mathrm d_pt$ is, by construction, an element of the dual space $T_p^*M$, hence a $(0,1)$ tensor; letting $p$
vary, a $(0,1)$ tensor \emph{field} $\mathrm dt$. Of a $(0,1)$ tensor field one can take the covariant derivative in any direction; the result is again a
$(0,1)$ tensor field (in coordinates one small index down). We require that it be zero everywhere, for every vector field $X$
\ts{8}{1}{12:08}: $\nabla_X\,\mathrm dt=0$. (In the lecture ``for every $X$ in $T_pM$''; corrected on the spot to ``for all vector fields $X$'', ``then you have disposed of it in one go''.)

These are, first of all, requirements, and they concern nothing but the relation of the objects of the quintuple among each other: for the
absolute time Z1 directly, and Z2 is a kind of compatibility condition with the covariant derivative \ts{8}{1}{13:32}. In the course of the lecture we shall see that
these two properties achieve, mathematically, exactly what the flowery formulation says \ts{8}{1}{14:13}.
So from now on, if somebody asks what ``the absolute time flows uniformly'' means, you say: the Newtonian spacetime is this quintuple, and by this we mean
$\nabla_X\,\mathrm dt=0$ for all vector fields $X$. If someone asks what it means that spacetime is foliated into three-dimensional spaces, you say:
that the $(0,1)$ tensor field $\mathrm dt$ defined by the absolute time is nowhere the zero map. Finished \ts{8}{1}{14:33}.
Are there questions on what stands there? None \ts{8}{1}{15:15}.

\begin{secondary}[On the word ``quintuple'']
The lecture says ``quintuple'' and then lists the five objects $M,\mathcal O,\mathcal A,\nabla,t$ (the manifold counts as three of them, so ``three ingredients''
above). The same count appeared for the affine manifold of \cref{def:cf-affine}, a quadruple.
\end{secondary}

\subsection{Five remarks}\label{ssec:nk-remarks}

There follow five remarks \ts{8}{1}{15:15}.

\paragraph{Remark 1: which connection, and which Newtonian spacetime?}
One can define the Newtonian spacetime in this way, a smooth manifold with a covariant derivative or linear connection, but one asks rightly: \emph{which}
connection do I choose? Also one might be bothered that the definition says ``\emph{a} Newtonian spacetime'' and not ``\emph{the}'': are there several?
There are, as we shall see \ts{8}{1}{15:54}. For today \emph{we do not commit}: we calculate with a general linear connection, some
$\Gamma$'s of which we know the properties; how they look concretely comes later. But a small outlook is allowed, so that one knows where the hare is running
\ts{8}{1}{16:14}.

\emph{Outlook A, hypothetical.} For the completely unrealistic case that there is \emph{no inhomogeneous gravitation}: a homogeneous gravitational field (the same everywhere) is what one
uses for the stone throw and cannonballs, but it is ``actually no gravitational field at all'' (a remark that can only be understood in the future; it will be shown).
An inhomogeneous one is for instance the field around a spherical mass, with arrows pointing towards it \ts{8}{1}{17:59}.
``No inhomogeneous gravitation'' here means that one may neglect the inhomogeneity on the scale one looks at; but we consider the \emph{whole}
spacetime at once, and if we are so megalomaniac, a purely homogeneous field is rather unrealistic. (Schuller offers the infinite flat plate: it has a field on both sides,
homogeneous on each, but pointing in opposite directions, hence globally not homogeneous either; ``I would not know how to produce a global homogeneous gravitational
field''. The transcript garbles part of this digression \ts{8}{1}{19:29}.) In that case it follows from the measurements described by the first Newtonian axiom
(which was a measuring prescription for what straight lines and uniform motion are) that the Newtonian spacetime (NRZ) is \textbf{flat}
\ts{8}{1}{20:33}. ``I stress this is an outlook, the first axiom will be made precise only today.''

What is dear to Schuller is to give a proper definition of flat \ts{8}{1}{21:17}.
\begin{definition}[New concepts: flat Newtonian spacetime (flach)]\label{def:nk-flat}
A Newtonian spacetime is called \textbf{flat} if there is a sub-atlas $\bar{\mathcal A}\subseteq\mathcal A$ of the atlas of the definition (a subset of the charts which still covers the whole
manifold, still called an atlas) such that for every chart $(U,x)\in\bar{\mathcal A}$ the connection coefficients vanish,
$\Gamma^a{}_{bc}{}_{(x)}=0$ \ts{8}{1}{21:38}.
\end{definition}
In slimmer words: one can always find, locally but covering the whole manifold, coordinates in which the Christoffel symbols vanish. Even then, if one goes to a chart not in
$\bar{\mathcal A}$ the $\Gamma$'s need no longer vanish (see \cref{sec:cf-connections}); but that such a collection of charts exists at all is a special property \ts{8}{1}{23:09}. In textbooks one often finds
the definition ``flat if the Riemann curvature tensor vanishes'' (met on the problem sheet): this is a necessary and sufficient condition for such an atlas $\bar{\mathcal A}$ to exist, but the
fundamental definition is the one above \ts{8}{1}{23:30}. It is a dream situation: one works always in $\bar{\mathcal A}$ with all $\Gamma=0$, which
simplifies a lot (leaving out the inhomogeneous gravitation, why this still works for homogeneous gravitation comes later). Only if one then goes to coordinates adapted to
the problem, e.g.\ polar coordinates, the nice property is lost again: there is a \emph{trade-off} between simple calculation on one side and the $\Gamma$'s on
the other. ``We simply must deal with the $\Gamma$'s, but that is the definition of flat'' \ts{8}{1}{24:12}.

\emph{Outlook B, complementary.} In the presence of an inhomogeneous gravitational field one must interpret it as a \emph{non-flat} spacetime \ts{8}{1}{25:15}.
Before continuing there should be a cry: ``wait, a gravitational field is a force field, gravitation is a force''. That is how it is learnt, and it is also right: one \emph{can} keep
spacetime flat and insist that gravitation is a force, which has nothing to do with this lecture and comes in the next. But recall the first Newtonian axiom in the old formulation:
a particle \emph{on which no force acts}. We discussed already that gravitation is universal, acts on every mass and cannot be shielded: with two, three, five or any number of particles in the
universe, gravitation always acts. If one describes it as a force, the first axiom, which speaks of particles on which no force acts, is \emph{deprived of its whole domain of validity}:
it would be void, since a force, at least gravity, acts on every particle \ts{8}{1}{26:20}. ``One can do it, but if you look at the first axiom, it is very brave.'' So we do
not: gravitation is not a force, it corresponds to a curvature of spacetime in Newtonian theory; how exactly, in one of the next lectures \ts{8}{1}{28:03}.
Hence: in the presence of an inhomogeneous gravitational field one must interpret this as a non-flat spacetime, if one does not want to deprive the first axiom of its
validity domain \ts{8}{1}{28:45}. And this would be fatal, because the first axiom is our \emph{measuring prescription for what is straight}. Recall the handful of sand in the spaceship, somewhere in
space with nothing to do: I scatter the sand, the grains fly along paths in what I perceive as space (to come), and the first axiom tells me, if the particles
have, for example, no charge coupling them to any fields (more in a moment), to identify these paths as the straight lines. If the first axiom had no domain of
validity left I could never do this \ts{8}{1}{29:18}. So one must treat gravitation this way; afterwards, for computing, he will show a trick to recast it again as a force, an equivalent
problem: ``only a computational vehicle, so that you can deal with gravity like with the rest; computationally it does not matter so much, but conceptually it is a hammer''
\ts{8}{1}{30:19}. This was an outlook, so that it is known where the road leads \ts{8}{1}{30:40}.

\begin{coursenote}[How to read the two outlooks]
Outlook A and B are announcements, not results of this lecture: the claim that homogeneous gravity is ``no gravitational field'' and that the inhomogeneous
one is curvature is derived only in \cref{ch:c11}. Note also that A and B are consistent with Z2 and flatness being independent:
the definition of a Newtonian spacetime above does \emph{not} require flatness. (Added; not said in class.)
\end{coursenote}

\paragraph{Remark 2: classification of vectors.} The second remark concerns the absolute time, which allows a classification, of great importance, of \emph{all vectors at a point}
$p\in M$ into three disjoint classes \ts{8}{1}{30:40}. Let $X\in T_pM$.
\begin{definition}[New concepts: spatial, future-pointing and past-pointing vectors]\label{def:nk-vectors}
A vector $X\in T_pM$ is called
\begin{itemize}
  \item \textbf{spatial} (r\"aumlich; the lecture first said ``raumartig'', then corrected) if $\mathrm d_pt(X)=X(t)=0$ \ts{8}{1}{31:26};
  \item \textbf{future-pointing} (zukunftsweisend) if $\mathrm d_pt(X)>0$ \ts{8}{1}{32:31};
  \item \textbf{past-pointing} (vergangenheitsweisend) if $\mathrm d_pt(X)<0$ \ts{8}{1}{32:51}.
\end{itemize}
\end{definition}
It is clear that this can be done: take any vector at $p$, let $\mathrm dt$ act at $p$, and out comes a real number, which is zero, positive or negative; so one knows to which species the vector
belongs \ts{8}{1}{33:34}. A picture comes soon, but first a few notions are needed \ts{8}{1}{33:55}.

\paragraph{Remark 3: events.} It has become customary, and is a rather nice picture, that every point $p\in M$ of spacetime has a special name: a \textbf{event} (Ereignis)
\ts{8}{1}{33:55}. Why give the thing another name? Roughly: a point in \emph{space} would be a place; a point in \emph{time} (on the time axis, which we do not have) would be a moment. But what is a point of
\emph{spacetime}? One is tempted to say ``a time and a place'', ``but very dangerous, that is written nowhere''. An event is, informally, a bomb explosion (hopefully not now and here)
\ts{8}{1}{34:58}. This is only terminology. The number $t(p)$ (a number: $t$ is a smooth function that eats a point) is called the \textbf{absolute time instant} of the event $p$
\ts{8}{1}{35:19}. This exists in Newton, not in Einstein: it is the time at which the event takes place \ts{8}{1}{36:00}.

\paragraph{Remark 4: the foliation into spaces.} The absolute time foliates the Newtonian spacetime completely into three-dimensional so-called \emph{spaces}
(``so far there was no talk of them'') \ts{8}{1}{36:21}.
\begin{definition}[New concepts: the space at absolute time $\tau$]\label{def:nk-space}
For $\tau\in\R$ the \textbf{space at the absolute time instant $\tau$} is the set of events
\[
  \mathcal R_\tau:=\{\,p\in M\mid t(p)=\tau\,\}\;=\;t^{-1}(\{\tau\})
\]
of all events whose absolute time instant is $\tau$ \ts{8}{1}{37:06}.
\end{definition}
The complete foliation is the statement that the union of all these spaces is the whole spacetime,
\begin{equation}\label{eq:nk-foliation}
  M\;=\;\dot{\bigcup_{\tau\in\R}}\;\mathcal R_\tau ,
\end{equation}
a \emph{disjoint} union: a point cannot lie in $\mathcal R_\tau$ for two different $\tau$, because every point has a definite absolute time \ts{8}{1}{37:50}. One may object: but what about the
absolute time standing still somewhere? That is excluded by Z1, $\mathrm dt$ is nowhere the zero map: time does not stand still anywhere \ts{8}{1}{38:11}.
On problem sheet~7 (the one handed out today was sheet~6) one proves that these spaces $\mathcal R_\tau$ are manifolds, of dimension~3; the instructions
are given there and one needs the property Z1 \ts{8}{1}{38:31}.

Now a picture \ts{8}{1}{39:16}. Spacetime cannot be drawn four-dimensional; so for the pictures the dimension is reduced by one: $\dim M=3$ (``which is actually not a
Newtonian spacetime'') and the spaces two-dimensional, everything else unchanged \ts{8}{1}{39:38}. Imagine the three-dimensional spacetime in perspective, with
the nowhere overlapping level surfaces of $t$: $t=1$, $t=3.14$, $t=-0.01$ below. The spacetime consists of such a foliation into spaces. Everything on one such surface is the space
$\mathcal R_{3.14}$, the next $\mathcal R_1$, $\mathcal R_{-0.01}$: ``roughly a photograph of the space at the corresponding instant, which the absolute time gives globally''
\ts{8}{1}{40:00}. The function $t$ can look all sorts of ways on the spacetime, but it looks roughly like this because $\mathrm dt$ is nowhere the zero map \ts{8}{1}{41:01}.

\paragraph{A student asks about the spaces drawn.}\label{par:nk-student} A student asks (inaudible, ``no, the spaces I draw\dots'') about the drawn spaces; Schuller: the spaces drawn are only
representatives of the \emph{continuum} of spaces in between: all of spacetime is the spacetime, and on it a curve and at each point a tangent vector; the drawn $\mathcal R_3$, $\mathcal R_{2.7}$, $\mathcal R_{-100}$ are only an orientation: in between lie also
$\mathcal R_{2.8}$, $\mathcal R_{2.9}$, uncountably many \ts{8}{1}{41:21}. \emph{Why disjoint?} Every point lies in exactly one of them (take a point, compute $\tau=t(p)$, then $p\in\mathcal R_{t(p)}$) and
taking all $\tau$ one gets every point \ts{8}{1}{41:42}. (The part of the transcript between 41:21 and 42:23 mixes the question and Schuller's answer; the reading of the
question is partly reconstructed [?].)

Now see how the whole thing is set up: spacetime is there \emph{primarily}. Then there is the absolute time $t$; only from spacetime and the absolute time do I define the absolute time instant of an event, and what the events at the same time are. ``We do not
start with spaces and glue them to a spacetime, we take the whole spacetime and identify in it what we call absolute time, the smooth function $t$, and what we call absolute space at instant $\tau$,
namely $\mathcal R_\tau$.'' It sounds simple, but it is of incredible importance that it is done this way round \ts{8}{1}{42:23}.

\paragraph{Remark 5: stratified atlases.} The fifth remark is, technically, \emph{very} important for calculations. It is by no means necessary, but unbelievably clever for simplifying calculations: to choose
a so-called \textbf{stratified atlas} \ts{8}{1}{43:03}. Look: on a Newtonian spacetime we have in particular the absolute time $t$, a function everywhere, smooth and therefore also continuous, $M\to\R$. Use this one function
as the first of the four chart component functions \ts{8}{1}{44:12}. In a chart $(U,x)$ one can always split the chart map into its components $x^0,x^1,x^2,x^3$; in spacetime we count, for
traditional reasons, not from $1$ to $4$ but from $0$ to $3$ (``Schuller wipes the blackboard again, intolerable''). And we require that the zeroth coordinate function be the restriction of the absolute time to $U$
\ts{8}{1}{45:39}.
\begin{definition}[New concepts: stratified atlas]\label{def:nk-stratified}
An atlas $\mathcal A$ of a Newtonian spacetime is called \textbf{stratified} (stratifiziert) if every chart $(U,x)\in\mathcal A$ has as zeroth coordinate function the restriction of the absolute time,
\[
  x^0\;=\;t|_U
\]
\ts{8}{1}{47:04}.
\end{definition}
A chart map must be continuous, going from the chart domain to the reals; we have on the whole Newtonian spacetime a global function fulfilling these conditions, hence it is ``incredibly clever'' to choose this zeroth coordinate function for every
chart of such an atlas \ts{8}{1}{46:21}. Does one have to? No, one can choose differently, but it brings nothing \ts{8}{1}{47:46}. \textbf{Convention} (tacit assumption, from now on): the atlas $\mathcal A$ occurring in the definition of a Newtonian spacetime
is a stratified atlas. ``This is an assumption without physical content'': nothing physical has been required, it is only for practicality, to save many lines in the computations
\ts{8}{1}{48:28}.

Two consequences. \emph{First}, for every chart $(U,x)$ of the stratified atlas the maps $(U\cap\mathcal R_\tau,\,x^1,x^2,x^3)$ are charts of the space $\mathcal R_\tau$ \ts{8}{1}{49:11}. ``Now you see why one starts counting at zero'': the $1$-,
$2$- and $3$-coordinate of a space $\mathcal R_\tau$ are already the $1$-, $2$- and $3$-component of the spacetime chart \ts{8}{1}{50:16}. Convention for indices from now on: Latin indices $a,b,\dots$ run from $0$ to $3$ (spacetime indices), Greek indices $\alpha,\beta,\dots$ only from $1$ to $3$
\ts{8}{1}{1:17:27}.

The picture (still for $\dim M=3$): in the \emph{chart image} we have the $x^0$-direction and the other directions ($x^1$, $x^2$ here). If one chooses a stratified atlas, each space $\mathcal R_\tau$ appears in the chart as one of these planes
($\mathcal R_0$, $\mathcal R_1$, \dots), hence the name ``stratified''. In the charts it looks more as one always imagined it: one takes pictures of space and lays them on top of each other, and has a picture of spacetime. And since one
can in fact draw only in the chart, in a stratified atlas the spaces lie so nicely as planes in a chart: ``simply a nice intuition'' \ts{8}{1}{50:58}.

\begin{nfigure}
\begin{tikzpicture}[scale=0.9]
  \foreach \z/\n in {0/{\mathcal R_{-1}},1.6/{\mathcal R_0},3.2/{\mathcal R_1}} {
    \draw[fill=blue!8,draw=blue!60!black] (0,\z) -- (4,\z) -- (5.4,\z+0.9) -- (1.4,\z+0.9) -- cycle;
    \node[blue!60!black,right] at (5.4,\z+0.5) {$\n$};
  }
  \draw[-{Stealth},thick,color=red!70!black] (2.0,1.95) -- (3.4,1.95) node[right] {spatial};
  \draw[-{Stealth},thick,color=green!50!black] (2.0,1.95) -- (2.0,3.0) node[above] {future-pointing};
  \draw[-{Stealth},thick,color=orange!80!black] (2.0,1.95) -- (2.0,0.9) node[below] {past-pointing};
  \fill (2.0,1.95) circle (1.5pt);
  \draw[-{Stealth}] (-0.4,0) -- (-0.4,3.9) node[above] {$x^0=t$};
\end{tikzpicture}
\caption*{Chart image of a stratified atlas for $\dim M=3$: the spaces $\mathcal R_\tau$ are the horizontal planes
(only representatives of the continuum between them), and the three kinds of tangent vector at an event.}
\end{nfigure}
Given this one can at once draw the difference between a spatial, a future-pointing and a past-pointing vector \ts{8}{1}{52:40}: at a point $p$, which lies in particular in a space, a \emph{spatial} vector is one that lies
tangentially to a space (easy to see, since it is annihilated by the gradient of the absolute time, whose level surface the space is); one pointing upwards is future-pointing, one pointing backwards
past-pointing \ts{8}{1}{53:01}. How does one tell? From where to where does the absolute time increase? ``This is completely encoded in the definition I gave''. Had the
spaces been labelled the other way ($\mathcal R_3$ at the bottom, $\mathcal R_1$ above), then the other would be future-pointing; it is the old \emph{gradient argument}: $\mathrm dt$ points in the direction of largest increase
(with the cosmetic flaw that it is a covector, not a vector, and hence does not ``point''; but roughly) \ts{8}{1}{53:45}.

\emph{Second consequence}, of the greatest importance \ts{8}{1}{54:50}: in a chart of a stratified atlas,
\[
  (\mathrm dt)_a\;=\;\frac{\partial x^0}{\partial x^a}\;=\;\delta^0{}_a ,\qquad a=0,1,2,3,
\]
since $\mathrm dt$ is identified with $\mathrm dx^0$ and $\partial x^0/\partial x^a$ is $1$ for $a=0$ and zero otherwise, at every point of spacetime
\ts{8}{1}{55:11}. Why is this so important? Because then Z2 becomes a condition on the $\Gamma$'s \ts{8}{1}{55:52}. Whenever ``chart'' is said from now on, a chart of a stratified atlas is meant
\ts{8}{1}{56:12}. In such a chart the condition Z2 reads, for the $a$-th component, with $X^m$ the components of an arbitrary vector field in this chart,
(``it is a covector, we know how to do it; there is a minus'') \ts{8}{1}{56:34}:
\[
  X^m\frac{\partial}{\partial x^m}(\mathrm dt)_a\;-\;\Gamma^s{}_{a\,m}\,(\mathrm dt)_s\,X^m
  \;=\;0-\Gamma^0{}_{a\,m}\,X^m\;=\;0 .
\]
The last index $m$ of $\Gamma$ goes with the vector field $X$ in whose direction one differentiates; $s$ is the dummy index summed with the components $(\mathrm dt)_s=\delta^0{}_s$ \ts{8}{1}{57:16}. The first term is the derivative of
constants, which is zero (one first turns it into a partial derivative with the chart, but nevertheless it is zero) \ts{8}{1}{57:58}. Hence $\Gamma^0{}_{am}X^m=0$ for every $X$, at every point, which means, independently of $X$:
\begin{equation}\label{eq:nk-Z2-comp}
  \Gamma^0{}_{ab}\;=\;0\qquad\text{for all }a,b\in\{0,1,2,3\}
\end{equation}
\ts{8}{1}{58:18}. This is the coordinate formulation of the condition Z2 in a stratified atlas: the $\Gamma$'s which have a $0$ as upper index (the zeroth component of the spacetime chart), \emph{whatever the lower indices}, vanish \ts{8}{1}{59:00}.
``You will see that this insight is quite enormous; we shall look at it right away'' \ts{8}{1}{59:21}.

\begin{coursenote}[Index placement in (\ref{eq:nk-Z2-comp})]
The condition comes out as $\Gamma^0{}_{am}X^m=0$, with the free index $a$ in the \emph{second to last} position and $m$ (contracted with $X$) last, as always in the convention of \cref{def:cf-covariant-tensor};
with torsion freeness \eqref{eq:nk-sym} the order of the two lower indices is immaterial anyway. In the lecture the same formula is written with the free index on a slightly different slot; the result \eqref{eq:nk-Z2-comp} is not affected.
\end{coursenote}

% ==================================================================
\section{Particles}\label{sec:nk-particles}

Second section, before the break in six minutes: \emph{particles}. In Newton one speaks of bodies (K\"orper); what Newton calls bodies we call particles (Teilchen). The question is: what, please, \emph{is} a particle?
\ts{8}{1}{59:42}
\begin{definition}[New concepts: particle (Teilchen), world line, mass, charges]\label{def:nk-particle}
A \textbf{particle} on a Newtonian spacetime is a tuple $(\gamma,m,c_1,\dots,c_L)$, with $L$ some natural number, consisting of
\begin{enumerate}
  \item a smooth curve $\gamma\colon\R\to M$ \emph{on the Newtonian spacetime} (``exclamation mark, not in space''), called the \textbf{world line} (Weltlinie) of the particle, such that
        all its tangent vectors $X_{\gamma,\gamma(\lambda)}$ are \textbf{future-pointing}, for all $\lambda\in\R$ \ts{8}{1}{1:00:08};
  \item a positive real number $m>0$, called the \textbf{inertial mass} (tr\"age Masse, Inertialmasse) of the particle \ts{8}{1}{1:05:28};
  \item real numbers $c_0,\dots,c_L$ (no restriction on the sign), called the \textbf{charges} of various kinds of the particle \ts{8}{1}{1:06:15}, which couple to corresponding fields, force fields, $F_0,\dots,F_L$
        (``in quotation marks'': to become clear in the next lecture, on dynamics) \ts{8}{1}{1:06:37}.
\end{enumerate}
\end{definition}
(The charges are indexed ``$c_1$ to $c_L$'' at first and ``$c_0$ to $c_L$'' a moment later in the lecture: a slip, as Schuller noticed, which does not matter for what follows [?].)

Remarks. The curve is remarkable because it lives on the whole spacetime \ts{8}{1}{1:00:50}, and also because \emph{the whole of $\R$} was chosen as the parameter interval. Not yet entirely clear, but we shall see today that this choice, together with the first
Newtonian axiom, is nothing but the statement that \textbf{particles are indestructible} (unzerst\"orbar): in Newtonian physics particles can neither be created nor annihilated; whoever wants that has to tinker here, and we do not
\ts{8}{1}{1:01:11}. The curve has the name \emph{world line of the particle}. The world line is also the \emph{history} of the particle: if you know the world line in spacetime you are like God: you
look at it and know what has always been and always will be with the particle (a point made later \ts{8}{1}{1:13:47}).

\emph{Future-pointing tangent vectors} means, visualised in a chart of a stratified atlas: a world line traversed in the direction of increasing $x^0$, whose tangent at every place points ``upwards''
\ts{8}{1}{1:02:57}. A world line that turns back, running into the past, ``is no particle at all: that is a particle that decides at some time to run into the past''; somewhere it has a spatial tangent vector, and elsewhere even past-pointing ones;
``nothing of that is allowed'' \ts{8}{1}{1:03:43}. A particle is a tuple in which the curve has \emph{all} tangent vectors future-pointing \ts{8}{1}{1:04:04}.

\emph{The spaces drawn} are only representatives of the continuum of spaces between them (as above, answering a student) \ts{8}{1}{1:04:26}.

\emph{Mass and charges.} Positivity of $m$ matters: massless particles do not exist in Newtonian mechanics. They appear in relativity, where we also want to get \ts{8}{1}{1:05:50}.
\begin{intuition}[The trick of the charges: how to say that no force acts]\label{int:nk-charges}
In the first law one had to say that a particle should have \emph{no force acting on it}. But how do I know that no force acts? I want to establish the straight lines with the first axiom; a force acts, as I
know, only because there is a deviation from the straight line; so I would have to know in advance how to detect that I have no force. This is the trick around it: later we shall formulate the first axiom not for
``particles on which no force acts'', but for \emph{particles all of whose charges are zero}. The force fields may well be there, but the particle does not feel them, it does not couple to them. For example, let $c_1$ be the electric charge: in the simplest case the particle couples to an electric field; but if its electric charge is $0$, a huge electric field may be
there and the particle simply will not feel it \ts{8}{1}{1:06:58}. So we only need to understand particles as such tuples: world line, mass and charges \ts{8}{1}{1:08:22}.
\end{intuition}

\begin{aside}[Break: a message about the exercise groups]
Schuller substituted a sick tutor in one exercise group last Thursday and was impressed by what he saw in learning performance and the resulting understanding of the definitions
(``fantastic; exactly as it should be: the kind of preparation one can only wish for''); he asks students to compare their group with it \ts{8}{1}{1:08:43}.
He asks for a quick show of hands: who is satisfied with their group, who is not (he did not look at who was unsatisfied); dissatisfied students should write him an e-mail, also anonymously, say what they do not like
and, best, how it could be improved, or come to talk in person; they will get an answer \ts{8}{1}{1:09:45}. Break until twenty past \ts{8}{1}{1:08:43}.
\end{aside}

% ==================================================================
\section{The first Newtonian law, precisely}\label{sec:nk-first-law}

Third: the \emph{precise formulation of the first Newtonian law} \ts{8}{1}{1:10:26}. It uses only the notions we have, and we stare at the formulation until it is clear why it really is the usual first Newtonian law
\ts{8}{1}{1:10:48}.
\begin{theorem}[First law of motion (erstes Bewegungsgesetz)]\label{ax:nk-first}
The world lines of particles with vanishing charges (all charges $=0$) are autoparallels (of the covariant derivative $\nabla$ of the Newtonian spacetime) \ts{8}{1}{1:11:30}.
\end{theorem}
(It is an \emph{axiom}, boxed on the board; the box type of the notes marks it as a statement; it is not proved but postulated.) ``Done.'' In contrast to the formulation of the first day, every word in this sentence is mathematically
crystal clear, and from now on this is our first axiom \ts{8}{1}{1:12:13}. What we do next is to look at it and massage it until we see that it is exactly what the usual formulation says, but now everything is clear
\ts{8}{1}{1:12:34}.

\subsection{Uniform motion: the affine parametrisation}\label{ssec:nk-affine}

\begin{theorem}[World lines are affinely parametrised by the absolute time]\label{thm:nk-affine}
The first law implies that the world lines of particles with vanishing charges are \textbf{affinely parametrised}: the curve parameter $\lambda$ differs from the absolute time $t(\gamma(\lambda))$ of the event $\gamma(\lambda)$ at most by an affine
transformation, once and for all chosen constants $a\in\R\setminus\{0\}$ and $b\in\R$,
\[
  t\bigl(\gamma(\lambda)\bigr)\;=\;a\,\lambda+b ,
\]
equivalently $\lambda=(t-b)/a$ \ts{8}{1}{1:12:56}.
\end{theorem}
(Said first as ``$\lambda=a\,t+b$''; the proof below gives the form above, where the constants are the reciprocals ($\lambda = (1/a)(t-b)$), as Schuller remarks at the end; the constants are renamed anyway.)
In words: ``the parameter of this world line is the history; one can also call it the history of the particle''. The parameter of the world line is distinguished from the absolute time of the event $\gamma(\lambda)$: ``you say: but the absolute time is a
\emph{function}, the parameter a \emph{number}; correct, but they differ at most by this'' \ts{8}{1}{1:14:33}. This is something special: a world line one may parametrise as one likes, run slowly, run first slowly then fast, whatever slowly and fast mean;
one can reparametrise and it is still the same world line, the same history of the particle. That is correct; but from the first law this follows, and the proof is fantastically simple \ts{8}{1}{1:15:42}.

\begin{proof}
From the first law we know that, in the chart representation of the world line in an arbitrary chart $(U,x)$ of the stratified atlas, with $q^a(\lambda):=x^a(\gamma(\lambda))$, the autoparallel equation holds,
\[
  \ddot q^a(\lambda)+\Gamma^a{}_{bc}\bigl(\gamma(\lambda)\bigr)\,\dot q^b(\lambda)\,\dot q^c(\lambda)=0 ,
\]
here with dots for derivatives with respect to the parameter $\lambda$ (\cref{thm:cf-autoparallel}; everything in this chart, the dependence on $x$ is not written) \ts{8}{1}{1:16:02}. Split off the zeroth equation from the three spatial ones:
\begin{align*}
  \ddot q^0+\Gamma^0{}_{bc}\,\dot q^b\dot q^c&=0 ,\\
  \ddot q^\alpha+\Gamma^\alpha{}_{bc}\,\dot q^b\dot q^c&=0\qquad(\alpha=1,2,3)
\end{align*}
\ts{8}{1}{1:17:07}, where the sum over $b,c$ runs over all spacetime indices ($0$ to $3$); together these are four equations, the chart representation of the autoparallel
condition on the world line \ts{8}{1}{1:18:09}. We look at the first one for now: by Z2 in the form \eqref{eq:nk-Z2-comp} we have $\Gamma^0{}_{bc}=0$ for all $b,c$, so the term drops out completely and
\[
  \ddot q^0(\lambda)=0
\]
in every stratified chart \ts{8}{1}{1:18:50}. This is easy to solve: $q^0(\lambda)=a\,\lambda+b$ for real numbers $a,b$, the only solutions of this equation \ts{8}{1}{1:19:32}. But what is $q^0$? It is
$x^0\circ\gamma$, the zeroth component of the curve, and in a stratified atlas the zeroth coordinate function is the absolute time \ts{8}{1}{1:20:15}: so $t(\gamma(\lambda))=a\lambda+b$. Solving for $\lambda$,
$\lambda=(t-b)/a$, and $a\ne0$ must hold: otherwise $t$ would be constant along the curve and the tangent vector could not be future-pointing (its action on $t$ would vanish) \ts{8}{1}{1:20:37}. This is the claim (with $a,b$ the constants above). \qedhere
\end{proof}
\begin{coursenote}[Why $a\neq0$, and the sign of $a$]
In the lecture ``$a$ must be nonzero, also fine'' is read off from the formula $\lambda=(t-b)/a$ which would make no sense; the reason in terms of the definition of a particle is the one in the proof
(the world line has future-pointing tangent vectors, $\mathrm dt(\dot\gamma)>0$). More precisely this even gives $a>0$: $\mathrm dt(\dot\gamma)=\dot q^0=a>0$. (Added; not said in class.)
\end{coursenote}

Consequence of this skirmish: because of the first law this initially completely freely choosable parameter is restricted to parameters that essentially run like the absolute time; one may choose another unit of time or move the offset
\ts{8}{1}{1:21:19}. In particular, if one differentiates a component of the curve once with respect to $\lambda$ one could just as well differentiate with respect to $t$, by the chain rule with $\lambda$ as a function of $t$:
$\mathrm dq^m/\mathrm dt=\dot q^m\,\mathrm d\lambda/\mathrm dt=\dot q^m/a$ (written with a prime $q'$ for $\mathrm d/\mathrm d\lambda$ and a dot for $\mathrm d/\mathrm dt$ in the lecture)
\ts{8}{1}{1:22:03}; differentiating twice, the factor is squared: $\mathrm d^2q^m/\mathrm dt^2=q''^m/a^2$ \ts{8}{1}{1:23:05}. In the autoparallel equations, each of the terms $q''$ gets a factor $a^{2}$ and each product $q'q'$ gets $a\cdot a$: the right-hand side is zero and the
factor $a^2$ can be thrown out. Hence one may replace $q'^a$ by the dotted derivative with respect to the absolute time, and it is still an autoparallel equation \ts{8}{1}{1:23:46}:
\begin{equation}\label{eq:nk-autopar-t}
  \ddot q^a(t)+\Gamma^a{}_{bc}\,\dot q^b(t)\,\dot q^c(t)\;=\;0 ,\qquad\text{with }\dot{}=\frac{\mathrm d}{\mathrm dt}.
\end{equation}
(A student corrects a missing double prime on the board, ``a double stroke has to stand there''; Schuller thanks him \ts{8}{1}{1:24:35}.) We have taken an autoparallel arbitrarily parametrised and the first law tells us
at once that this parameter must be affine, linear with an offset in the absolute time. In words, this is exactly the \textbf{uniformity of the motion} (Gleichf\"ormigkeit der Bewegung) mentioned in the traditional formulation of the first law
\ts{8}{1}{1:24:58}. \emph{Not} the straightness, but the uniformity. And it has nothing at all to do with inertial systems: the traditional formulation says that ``in an inertial system a particle moves
rectilinearly and uniformly in space''; ``nope: in \emph{any} system, says our first axiom'' \ts{8}{1}{1:25:43}.

\begin{coursenote}[Which charts?]
The proof works in every chart of a \emph{stratified} atlas (the only place where it matters is $x^0=t|_U$ and $\Gamma^0{}_{ab}=0$); uniformity thus holds in every such chart, inertial or not. Charts that are not stratified
(e.g.\ rotating coordinates that mix $x^0$ with space) are not covered by the argument, but the statement $t\circ\gamma=a\lambda+b$ is chart independent. (Added; not said in class.)
\end{coursenote}

% ==================================================================
\section{Inertial systems and straight-line motion}\label{sec:nk-inertial}

To speak of the \emph{straightness} of the motion one defines inertial systems \ts{8}{1}{1:26:29}. This is one of those notions of which the literature always refers to the fixed-star sky and says all sorts of philosophical things;
``we can of course also define it precisely''.

\begin{definition}[New concepts: global inertial system (globales Inertialsystem)]\label{def:nk-inertial}
A \textbf{global inertial system} of a Newtonian spacetime is an atlas $\mathcal A_{IS}\subseteq\mathcal A$, formed entirely of charts of the stratified atlas, such that for every chart $(U,x)\in\mathcal A_{IS}$
\begin{equation}\label{eq:nk-inertial}
  \nabla_{\partial/\partial x^0}\,\frac{\partial}{\partial x^m}\;=\;0\qquad\text{for }m=0,1,2,3
\end{equation}
\ts{8}{1}{1:26:50}. (Here $\partial/\partial x^0$ is the vector field of the chart; one could write $\partial/\partial t$ as well, ``the same''.)
These are four vector equations. For $m=0$ the condition is called the \textbf{non-acceleration condition} (Nichtbeschleunigungsbedingung), for $m=1,2,3$ the \textbf{non-rotation conditions} (Nichtrotationsbedingungen).
(``I should call them conditions rather than equations, as I put them on the charts which lie in such an atlas.'') \ts{8}{1}{1:28:20}
\end{definition}
In the literature, in particular in general relativity, this is found under the name \emph{Fermi--Walker transport}; one need not look it up, ``we are doing it here'' (as said in class, not checked) \ts{8}{1}{1:29:48}.

The first and probably most important remark: a global inertial system \emph{may or may not exist} in a given Newtonian spacetime; there is no guarantee that there is one. Example: in the presence of an inhomogeneous gravitational field there is none
\ts{8}{1}{1:30:09}. A second remark was dropped for lack of time and put on a problem sheet \ts{8}{1}{1:31:17}.

\subsection{The inertial conditions in components}\label{ssec:nk-inertial-comp}

Now we evaluate the inertial-system conditions in the components of a chart (suppose there is one; ``we assume, we must compute it, one can see it right away'') \ts{8}{1}{1:31:38}.
What do the conditions mean for our $\Gamma$'s? Start with $\nabla_{\partial/\partial x^0}\,\partial/\partial x^0$, the non-acceleration condition. The result is a vector field, with components $a$ up. By the rule of the covariant derivative
\begin{align*}
  \bigl(\nabla_{\partial_0}\partial_0\bigr)^a
  &=\frac{\partial}{\partial x^0}\,\delta^a{}_0\;+\;\Gamma^a{}_{mn}\,\delta^m{}_0\,\delta^n{}_0
  \;=\;0+\Gamma^a{}_{00},
\end{align*}
because the basis vector $\partial_0$ has the components $\delta^a{}_0$ in its own basis (a one in the zeroth slot, zero otherwise; hence the derivative term vanishes), the field being differentiated enters $\Gamma$ with the second to last
index, the direction with the last \ts{8}{1}{1:32:00}. The requirement is that this vanish: $\Gamma^a{}_{00}=0$. But we know already from Z2 that $\Gamma^0{}_{ab}=0$; so the new condition, not following from Z2,
is
\begin{equation}\label{eq:nk-noaccel}
  \Gamma^\alpha{}_{00}\;=\;0\qquad(\alpha=1,2,3)
\end{equation}
\ts{8}{1}{1:33:41}. In terms of the $\Gamma$'s the non-acceleration condition is the vanishing of these components: ``we started with completely general $\Gamma$, but now with a stratified atlas, and then an inertial system, they all drop like flies'' \ts{8}{1}{1:34:03}.
(A student points out an index which is $\alpha$ where $a$ was written. Schuller: for $a=0$ the equation reproduces $\Gamma^0{}_{00}=0$, which was already known from Z2, so the \emph{new} conditions are only $\Gamma^1{}_{00},\Gamma^2{}_{00},\Gamma^3{}_{00}$ \ts{8}{1}{1:34:32};
``I could write it for all $a$, which would be right, but I want to write economically what exactly it adds that is new''.) Question clear? Answer clear? \ts{8}{1}{1:35:35}.

Now the non-rotation conditions, $m\ne0$, with Greek $\mu=1,2,3$: $\nabla_{\partial_0}\partial_\mu$, $a$-th component \ts{8}{1}{1:35:57}:
\[
  \bigl(\nabla_{\partial_0}\partial_\mu\bigr)^a\;=\;\frac{\partial}{\partial x^0}\,\delta^a{}_\mu+\Gamma^a{}_{mn}\,\delta^m{}_\mu\,\delta^n{}_0
  \;=\;0+\Gamma^a{}_{\mu0} ,
\]
the derivative term again zero (components constants) \ts{8}{1}{1:36:17}. The requirement of non-rotation is $\Gamma^a{}_{\mu0}=0$. For $a=0$ this is already known; but because of the torsion freeness \eqref{eq:nk-sym} $\Gamma^a{}_{\mu0}=\Gamma^a{}_{0\mu}$, so
if I have found one I have the other (``this is why we asked for torsion freeness at the beginning''). The new conditions are
\begin{equation}\label{eq:nk-norot}
  \Gamma^\alpha{}_{0\mu}\;=\;\Gamma^\alpha{}_{\mu0}\;=\;0\qquad(\alpha,\mu=1,2,3)
\end{equation}
\ts{8}{1}{1:36:59}. (``Greek indices run from 1 to 3''.)
So how was it with the global inertial system? It may not exist: if a spacetime has no atlas in which these $\Gamma$'s vanish. This is for instance the case with an inhomogeneous gravitational field: in it one can never find an atlas such that these
vanish; ``why, I shall show you in a next lecture'' \ts{8}{1}{1:38:09}. So that is the condition on the $\Gamma$'s from the inertial-system requirement, which one possibly cannot fulfil at all \ts{8}{1}{1:38:30}.

\subsection{What it brings: straight-line motion in space}\label{ssec:nk-straight}

What does all this bring? A wonderful insight \ts{8}{1}{1:38:51}. Recall that of the four autoparallel equations of the world line we looked above only at the zeroth. Now look at the \emph{spatial} ones, which also all hold \ts{8}{1}{1:39:13}. In an inertial system, with the dot for the
derivative with respect to the absolute time \eqref{eq:nk-autopar-t}:
\[
  \ddot q^\alpha+\Gamma^\alpha{}_{mn}\,\dot q^m\dot q^n=0 .
\]
The summation over $m,n$ runs over all spacetime indices; split it by cases \ts{8}{1}{1:39:54}:
\[
  \ddot q^\alpha
  +\underbrace{\Gamma^\alpha{}_{00}}_{0}\,\dot q^0\dot q^0
  +\underbrace{\Gamma^\alpha{}_{0\gamma}}_{0}\,\dot q^0\dot q^\gamma
  +\underbrace{\Gamma^\alpha{}_{\beta0}}_{0}\,\dot q^\beta\dot q^0
  +\Gamma^\alpha{}_{\beta\gamma}\,\dot q^\beta\dot q^\gamma=0 ,
\]
where in the last term the sum is over $\beta,\gamma=1,2,3$ only \ts{8}{1}{1:40:56}. (Same line as before, only the spacetime sum split up \ts{8}{1}{1:41:17}.) Now in an inertial system the first of these vanishes by non-acceleration \eqref{eq:nk-noaccel},
the second and third by non-rotation \eqref{eq:nk-norot} \ts{8}{1}{1:41:38}: ``an inertial system is not only one that is not accelerated in some direction, it is also one that does not rotate; you could be unaccelerated and yet rotate'', e.g.\ a rocket at launch is accelerated and
also rotates, ``certainly not an inertial system'' \ts{8}{1}{1:41:59}. What remains is
\begin{equation}\label{eq:nk-space-autopar}
  \ddot q^\alpha(t)+\Gamma^\alpha{}_{\beta\gamma}\bigl(q(t)\bigr)\,\dot q^\beta(t)\,\dot q^\gamma(t)\;=\;0 ,
\end{equation}
holding for spatial indices only \ts{8}{1}{1:42:43}. What kind of equation is this? It has a very definite name, we know it already: the \emph{autoparallel equation}, but not the one with which we started in spacetime: it is the
autoparallel equation \emph{in space} \ts{8}{1}{1:43:24}. The first Newtonian axiom speaks of an autoparallel in spacetime; from this we took a lot of information, and in an inertial system, since we are in a stratified atlas, the motion in a space $\mathcal R_\tau$ obeys the autoparallel equation in space
\ts{8}{1}{1:43:45}.

\begin{theorem}[First law, traditional form]\label{thm:nk-first-traditional}
In a global inertial system (using the first law \ref{ax:nk-first}), a particle with vanishing charges moves \textbf{uniformly} and \textbf{straight in space}: the first law of the traditional formulation \ts{8}{1}{1:44:25}.
\end{theorem}
Where did we derive it? We took the autoparallel equation in spacetime; of this we had already established that on the one hand it gives the \emph{uniformity} of the motion (above), and now it gives even the \emph{straightness} of the motion in space,
of which the traditional formulation of the first law speaks \ts{8}{1}{1:44:46}. So: from the crystal-clear formulation of the first Newtonian axiom, that the world lines (histories) of particles with vanishing charges are autoparallels in the
Newtonian spacetime, we have derived the full traditional first axiom: particles move in an inertial system uniformly and straight, if no force acts on them \ts{8}{1}{1:45:08}. What has happened to ``no force acts''? That is
the vanishing charges: when all charges are zero (to be shown next time), forces can be there, but the particle does not couple to them, it does not feel them \ts{8}{1}{1:45:49}.
\emph{Gravitation is not a force}; how it curves spacetime we shall see; but only by not treating gravitation as a force do we have an application domain of our first axiom at all. Next time: the deviation from uniform
straight motion, which leads to force fields \ts{8}{1}{1:46:10}. Thank you.

\begin{coursenote}[The lecture's chain of reasoning, summarised]
(i) $\Gamma^0{}_{ab}=0$ (Z2 in a stratified chart) gives uniformity: $t\circ\gamma$ affine in the parameter. (ii) The inertial conditions $\Gamma^\alpha{}_{00}=\Gamma^\alpha{}_{0\mu}=0$ then remove all mixed terms from the spatial autoparallel
equations, leaving the autoparallel equation of the spatial part $\Gamma^\alpha{}_{\beta\gamma}$. In a flat spacetime with all $\Gamma=0$ in the charts of $\bar{\mathcal A}$ (\cref{def:nk-flat}) this is $\ddot q^\alpha=0$: straight lines in the usual sense. In an inertial
system of a non-flat one the spatial $\Gamma^\alpha{}_{\beta\gamma}$ may remain. (Added; not said in class.)
\end{coursenote}

% ==================================================================
\section*{Key concepts}
\begin{keyconcepts}
\textbf{Retrieval questions.}
\begin{enumerate}
  \item Define a Newtonian spacetime: which objects, which dimension, which two conditions on the covariant derivative and the absolute time?
  \item What is the torsion tensor and what does torsion freeness say about the $\Gamma$'s?
  \item State Z1 and Z2; what do they mean in words and what do they give in a stratified chart ($\Gamma^0{}_{ab}=0$)?
  \item When is a Newtonian spacetime flat? How does this relate to the vanishing of the curvature tensor?
  \item Define spatial, future-pointing and past-pointing vectors, events, the absolute time instant and the spaces $\mathcal R_\tau$; why is the union disjoint and complete?
  \item What is a stratified atlas and why is it physically empty but computationally useful?
  \item Define a particle. Why is the interval of the world line all of $\R$ (indestructibility), and why $m>0$?
  \item State the first law in the precise form; prove that world lines of particles with vanishing charges are affinely parametrised by the absolute time.
  \item Define a global inertial system, with the acceleration and rotation conditions; write them in components; how does the first law then give straight, uniform motion in space?
  \item Why does one introduce charges instead of ``no force acts''? Why is gravitation not a force here?
\end{enumerate}
\textbf{Concepts.}
\begin{itemize}
  \item Newtonian spacetime, torsion freeness, absolute time Z1, Z2 (\cref{def:nk-spacetime}); flatness (\cref{def:nk-flat}).
  \item Spatial / future / past vectors, events, spaces $\mathcal R_\tau$, foliation (\cref{def:nk-vectors,def:nk-space}, \eqref{eq:nk-foliation}); stratified atlas and $\Gamma^0{}_{ab}=0$ (\cref{def:nk-stratified}, \eqref{eq:nk-Z2-comp}).
  \item Particles: world line, mass, charges (\cref{def:nk-particle}); first law and uniformity (\cref{ax:nk-first,thm:nk-affine}).
  \item Global inertial system and its conditions (\cref{def:nk-inertial}, \eqref{eq:nk-noaccel}, \eqref{eq:nk-norot}); straight motion in space (\eqref{eq:nk-space-autopar}, \cref{thm:nk-first-traditional}).
\end{itemize}
\end{keyconcepts}

\flushside
\end{paracol}
